\documentclass{article}
\usepackage{graphicx} % Required for inserting images
\usepackage{amsmath}
\title{Lesson 6 Challenge - Calculus}
\author{William Wang}
\date{September 2026}

\begin{document}

\maketitle
Since $f$ is continuous at $c$, by definition, for every $\varepsilon > 0$ there exists $\delta > 0$ such that whenever $x$ is in the domain and $|x-c| < \delta$,
\[
|f(x) - f(c)| < \varepsilon.
\]

Since $f(c) \neq 0$, choose
\[
\varepsilon = \frac{|f(c)|}{2} > 0.
\]

By continuity, there is a $\delta > 0$ such that for all $x \in (a,b)$ with $|x - c| < \delta$,
\[
|f(x) - f(c)| < \frac{|f(c)|}{2}.
\]

\textbf{Case 1: $f(c) > 0$.}

Then the inequality above gives
\[
-\frac{f(c)}{2} < f(x) - f(c) < \frac{f(c)}{2},
\]
\[
\frac{f(c)}{2} < f(x) < \frac{3f(c)}{2}.
\]
We have $f(x) > \frac{f(c)}{2} > 0$ for all $x$ in $(c-\delta, c+\delta)$. So $f(x)$ has the same sign as $f(c)$ throughout this interval.

\textbf{Case 2: $f(c) < 0$.}

Then $|f(c)| = -f(c)$, and the same inequality gives
\[
-\frac{-f(c)}{2} < f(x) - f(c) < \frac{-f(c)}{2},
\]
\[
\frac{3f(c)}{2} < f(x) < \frac{f(c)}{2}.\]
We have, $f(x) < \frac{f(c)}{2} < 0$ for all $x$ in $(c - \delta, c+\delta)$. Again $f(x)$ has the same sign as $f(c)$.

In either case, there is an interval $(c-\delta, c+\delta)$ about $c$ throughout which $f(x)$ has the same sign as $f(c)$.

\end{document}
